\Askhsh{14760}{4}
\begin{ylh}{2}
    \item 3.1. Εξισώσεις 1ου Βαθμού
    \item 4.2. Ανισώσεις 2ου Βαθμού
    \item 6.1. Η Έννοια της Συνάρτησης
    \item 6.2. Γραφική Παράσταση Συνάρτησης
\end{ylh}
Δίνεται η συνάρτηση: $g(x) = \left\lbrack \dfrac{1}{\sqrt[3]{x^{2} - x - 12}} \right\rbrack^{3} \cdot (x^{2} - 16)$.

\begin{alist}
\item Να βρείτε το πεδίο ορισμού της συνάρτησης $g$.
\monades{8}
\item Να δείξετε ότι $g(x) = \dfrac{x + 4}{x + 3}$ για κάθε $x$ στο πεδίο ορισμού της.
\monades{9}
\item Να βρείτε, αν υπάρχουν, τα σημεία τομής της γραφικής παράστασης της συνάρτησης $g$ με τους άξονες.
\monades{8}
\end{alist}

